\section{HOMOLOGY}\label{homology}
%The material in this section presents some not entirely trivial 
 %typesetting problems.  The source code presented here displays some
 %of the solutions Eduardo found.  At different times during the
 %writing similar problems were solved differently.  You should try to 
 %understand why things have been coded as they are in this section
 %and you should try to find your own optimal solutions to your
 %typesetting problems.
%
%
%The next line produces an indented paragraph to start the document
 %unit.  The LaTeX defaults start most units without indentations.
\hspace{\parindent}
We now introduce the concepts of simplicial module and of the homology
associated to a simplicial module.
As a variation of this, we recall the definition of cyclic module and its
cyclic homology first given by Loday \cite{Lodaybook}.

The letter $k$ will always denote a commutative ring with identity.

\subsection{Simplicial Homology of a Simplicial Module}
%The next line produces an indented paragraph to start the document
 %unit.  The LaTeX defaults start most units without indentations.
\hspace{\parindent}
We define simplicial modules and their homology. Also,
two examples are given.
The first example is the free simplicial module associated to a simplicial
set.
Our second example is the Hochschild complex of a unitary algebra $A$ over
a ring $k$. 

\begin{defn}
A simplicial $k$-module $M_*$ is a family of $k$-modules $\{ M_n \}_{n\geq 0}$
together with $k$-module homomophisms
$d_i \colon M_n\rightarrow M_{n-1}$
and
$s_i \colon M_n\rightarrow M_{n+1}$ 
for  $ 0\leq i\leq n$
satisfying the following identities:
\begin{align}
d_id_j&=d_{j-1}d_i \quad \text{for} \quad i<j. \label{eq:faces} 
\\
s_is_j&=s_{j+1}s_i \quad \text{for}\quad i\leq j \label{eq:degeneracies}
\\
d_is_j&=\begin{cases}
        s_{j-1}d_i & \text{if $i< j$}, \\
        id         & \text{if $i=j$ or $i=j+1$}, \\
        s_jd_{i-1} & \text{if $i>j+1$}.
        \end{cases} \label{eq:mix}
\end{align}
\end{defn}

In order to define the homology of a simplicial module define
$b \colon  M_n \rightarrow M_{n-1}$ to be the boundary map such that
\begin{equation*}
               b=\sum_{i=0}^n(-1)^id_i
\end{equation*}

Relations \eqref{eq:faces} imply $b^2=0$, so we obtain a
chain complex $(M_*,b)$ and we have the following definition:

\begin{defn}
Let $M_*$ be a simplicial $k$-module. The homology of the simplicial module
$M_*$ is given by
\[ H_n(M_*)=H_n(M_*,b) \]
where $H_n(M_*,b)$ is the $n$-th homology of the complex $(M_*,b)$.
\end{defn}
\begin{example}
A simplicial set $X$ is a family of sets $\{X_n\}_{n\geq 0}$ together with
functions $d_i$ and $s_i$ satisfying relations 
\eqref{eq:faces}-\eqref{eq:mix}.
A simplicial set gives rise to a free simplicial $k$-module denoted
$k[X]$ by taking $k[X]_n=k[X_n]$, the free $k$-module with basis
$X_n$. In this case we usually write
$H_*(k[X])=H_*(X;k)$.
Let $Y$ be a topological space and let $X=\Sin(Y)$ be the singular
complex of $Y$. This has 
\[ 
X_n=   \{ \sigma  \colon  
             \Delta^n\rightarrow Y : \text{$\sigma$ is continuous} \} 
\]
with faces and degeneracies induced by the cofaces and codegeneracies
of $\Delta^n$, where $\Delta^n$ is the geometric simplex of dimension $n$.
Then $H_*(X;k)$ is singular homology of the space $Y$ with coefficients in $k$.
\end{example}

\begin{example}\label{tensor_algebra}
Let $A$ be an associative $k$-algebra and let $A_*$ be the simplicial 
$k$-module defined by
\[ 
    A_n=A^{\otimes (n+1)}
       = A\otimes_k A\otimes_k \cdots \otimes_k A 
        \qquad \text{$(n+1)$ factors} 
\]
A generator of $A^{\otimes(n+1)}$ is denoted by $(a_0,\dots a_n)$.
Let
\[ 
d_i \colon A_n\rightarrow A_{n-1} 
\quad \text{and} \quad
   s_i \colon A_n \rightarrow A_{n+1} 
\]
be such that
\[ 
d_i(a_0,\dots ,a_n)=\begin{cases}
                (a_0,\dots , a_ia_{i+1},\dots a_n) & \text{if $0\leq i <n$}, \\
                (a_na_0,a_1,\dots ,a_{n-1})        & \text{if $i=n$}.
                \end{cases}
\]
\[ 
s_i(a_0,\dots ,a_n)=(a_0,\dots , a_i,1, a_{i+1},\dots ,a_n) 
            \quad \text{for} \quad 0\leq i\leq n
\]
It can be verified that these functions satisfy the simplicial identities
and, therefore, the family of modules $A_n=A^{\otimes (n+1)}$ together with
these maps $d_i$ and $s_i$ is a simplicial module.
Its homology is known as the
Hochschild homology of the algebra $A$ and is denoted $HH_*(A)$.
\end{example}

\subsection{Cyclic Homology of a Cyclic Module}
%The next line produces an indented paragraph to start the document
 %unit.  The LaTeX defaults start most units without indentations.
\hspace{\parindent}
In this section we describe the definition of cyclic homology given by
Loday and Quillen \cite{Lodaybook}.
\begin{defn}
A cyclic module $M$ is a simplicial module together with an action of
the cyclic group
$C_{n+1}=\langle t_{n+1} : (t_{n+1})^{n+1}=1 \rangle$ 
on $M_n$, for each $n$, such that the following additional relations hold
\begin{equation}\label{cyclic_relations}
\begin{array}{lclcl}
d_it_{n+1}=t_{n}d_{i-1} 
        &  \text{and} 
              & s_it_{n+1}=t_{n+2}s_{i-1}     
                        &\text{for} 
                             & 1\leq i\leq n, 
\\
d_0t_{n+1}=d_n            
        & \text{and} 
              &  s_0t_{n+1}=t_{n+2}^2s_n 
                        &           
                              &
\end{array}
\end{equation}
where $d_i \colon M_n\rightarrow M_{n-1}$   
and $s_i \colon M_n \rightarrow M_{n+1}$ for $i=0,1,\dots , n$
are the simplicial structure maps.
\end{defn}

Next we will define the cyclic homology of a cyclic module.
The bicomplex we describe now will reappear in a more conceptual way, after
we develop $\Tor$ functors associated to cyclic and cocyclic modules.

If we define two operators $T_{n+1}$ and $N_{n+1}$ by
\begin{equation*}
T_{n+1}=(-1)^nt_{n+1} \quad \text{and} \quad N_{n+1}=1+T_{n+1}+\dots +T^n_{n+1}
\end{equation*}
then we have the following lemma:
\begin{lemma}\label{cyclic_1}
Let ${\cal C}(M)$ be the diagram of modules and homomorphisms
\begin{equation}\label{cyclic_bicomplex_1}
\xy \xymatrix{
\vdots \ar[d]_-b   
          &      \vdots \ar[d]_-{b'}      
                & \vdots \ar[d]_-b   
                      &  \vdots   \ar[d]_-{b'}        
                                 &                 
\\
M_2\ar[d]_-b 
          & \ar[l]_-{1-T_3}\ar[d]_-{b'} M_2 
                 & \ar[l]_-{N_3} \ar[d]_-{b}M_2 
                      & \ar[l]_-{1-T_3}\ar[d]_-{b'} M_2 
                                 & \ar[l]_-{N_3} \cdots       
\\
M_1\ar[d]_-b 
          & \ar[l]_-{1-T_2}\ar[d]_-{b'} M_1 
                 & \ar[l]_-{N_2} \ar[d]_-{b}M_1 
                       & \ar[l]_-{1-T_2}\ar[d]_-{b'} M_1 
                                 & \ar[l]_-{N_2} \cdots       
\\
M_0          
          & \ar[l]_-{1-T_1}             M_0 
                 & \ar[l]_-{N_1}            M_0 
                       & \ar[l]_-{1-T_1}             M_0 
                                  & \ar[l]_-{N_1} \cdots
}
\endxy
\end{equation}
where $b,b' \colon M_n \rightarrow M_{n-1}$ are defined by
\begin{equation*}
b=\sum_{i=0}^{n}(-1)^id_i \quad \text{and} \quad b'=\sum_{i=0}^{n-1}(-1)^id_i
\end{equation*}
Then the squares anticommute, the rows and columns
are complexes, so ${\cal C}(M)$ is a bicomplex.
\end{lemma}
\begin{proof}See \cite[page 52]{Lodaybook}
\end{proof}
\begin{defn}
Let $M_*$ be a cyclic module. 
The cyclic homology of $M_*$ is denoted $HC_*(M_*)$ and is defined as
\begin{equation*}
        HC_*(M_*)=H_*(\Tot({\cal C}(M)))
\end{equation*}
where ${\cal C}(M)$ is the the bicomplex of lemma \ref{cyclic_1}.
\end{defn}
\begin{example}
Let $A_*$ be the simplicial module defined in example \ref{tensor_algebra}. 
The action of the cyclic
group $C_{n+1}=\langle t_n : (t_{n+1})^{n+1}=1 \rangle $ on $A_n$ is given by:
\[ t_{n+1}(a_0,\dots ,a_n)=(a_n, a_1, \dots , a_{n-1}) \]
It can be verified that under these definitions $A_*$ becomes a cyclic module.
In this case we write 
$HC_*(A)$ instead of $HC_*(A_*)$ for the cyclic homology of $A$.
\end{example}

The vertical complexes $(M_*,b')$ in the bicomplex \eqref{cyclic_bicomplex_1}
are contractible, with contracting homotopy 
$h=T_{n+1}s_n \colon M_n \rightarrow M_{n+1}$.
Loday's Killing Contractible Complexes Lemma 
\cite[page 55]{Lodaybook} allows us to delete
all of these columns of the bicomplex 
\eqref{cyclic_bicomplex_1} obtaining the following
double complex:
\begin{equation*}
\xy \xymatrix{
\ar[d]_-b\vdots   & \ar[d]_-b\vdots   &  \ar[d]_-b \vdots \\
M_2 \ar[d]_-b  & \ar[d]_-b\ar[l]_-B M_1  &  \ar[l]_-B M_0  \\
M_1 \ar[d]_-b  & \ar[l]_-B M_0  &     \\
M_0            &                &
}
\endxy
\end{equation*}
Here 
$B=(-1)^{n+1}(1-t_{n+1})hN  \colon  M_n\rightarrow M_{n+1}$ 
and $b$ is the Hochschild boundary map. 
This bicomplex is called the ${\mathcal B}M$ complex 
or $(B\text{-}b)$-complex. 
For the particular case of the cyclic module associated to the algebra 
$A=T(V)$ as described in
\ref{tensor_algebra}, the $(B\text{-}b)$-complex
is used to give an explicit computation \cite{L-Q}. 
The result is given in the following theorem.
\begin{theorem}\label{tensor_cyclic_thm}
Let $V$ be a module over $k$, and let $T(V)$ be its tensor algebra. 
Then the cyclic homology
$HC_*(T(V))$ is given by:
\begin{equation*}
HC_*(T(V))=HC_*(BS^1,k)\oplus \sum_{m\geq 1} H_*(C_m,V^{\otimes m}) 
\end{equation*} 
where the generator $t_m$ of $C_m$ acts on $V^{\otimes m}$ as
\begin{equation*}
t_m(v_1v_2\cdots v_m)=v_mv_1\cdots v_{m-1}
\end{equation*}
\end{theorem}
\begin{proof}See \cite[page 581]{L-Q}.
\end{proof}
%This is the end of chp1a.tex
